\chapter{Não-dualismo}

\lettrine{A}{\space contribuição significativa} dos ensinamentos budistas reside
no que chamamos de não-dualismo. É a abordagem `nem isto nem aquilo' às
questões filosóficas. A religião monista tende a falar sobre o Um, o Deus Único,
o Todo, a Natureza Búdica ou a Mente Única, e isso é muito inspirador.
Recorremos às doutrinas monistas em busca de inspiração. Mas a inspiração é
apenas um nível da experiência religiosa, e é preciso ultrapassá-la. É preciso
abandonar o desejo de inspiração, ou a crença em Deus, na Unidade, na Mente
Única, na benevolência que tudo abrange ou na justiça universal.

Também não te estou a pedir que deixes de acreditar nessas coisas. Mas a prática
não dualista é uma forma de nos libertarmos de tudo isso, de percebermos o apego
a pontos de vista, opiniões e percepções, porque a percepção da nossa própria
mente é uma percepção, não é verdade? A percepção de uma benevolência universal
é uma percepção à qual nos podemos apegar. A Natureza Búddhica é uma percepção.
O Buddha é uma percepção. O Deus Único e tudo como sendo um sistema universal,
uma aldeia global, tudo é um e um é tudo e tudo é justo e tudo é bondoso, Deus
ama-nos: estas são percepções que podem ser muito agradáveis, mas continuam a
ser percepções que surgem e desaparecem. As percepções das doutrinas monistas
surgem e cessam.

Agora, o que é que isso faz, como experiência prática, quando deixamos as coisas
irem e elas cessam? O que resta, qual é o resíduo? É para isto que o Buddha
aponta ao ensinar sobre o surgimento e a cessação das condições.

Quando a percepção do eu cessa e todas as doutrinas, todos os ensinamentos
inspirados, todos os ditos sábios cessam, ainda há quem conhece a cessação. E
isso deixa-nos com uma mente em branco. O que há para agarrar? Então surge o
desejo de saber, de ter algo para agarrar. Às vezes podemos ver um pânico nas
nossas mentes: temos de acreditar em algo! «Fala-me da benevolência universal!»
Mas isso é o medo e o desejo a operar novamente, não é? «Quero acreditar em
algo! Preciso de algo em que acreditar! Quero saber que está tudo bem. Quero
apegar-me e acreditar nas percepções de unidade e totalidade.» E assim,
continua a existir esse desejo a operar, ao qual talvez não dês por isso e ao
qual ainda possas estar apegado. É por isso que a experiência religiosa é uma
experiência de desespero.

Na história da crucificação, a afirmação mais impressionante que Cristo fez foi:
«Pai, Pai, por que me abandonaste?» O que aconteceu àquele Pai que estava a
proteger Jesus? Até Deus o abandonou. É um grito angustiado, não é? A única
percepção com que «eu» podia contar dissolveu-se subitamente na sua mente. E
depois disso, claro, vem a aceitação e, em seguida, a Ressurreição, o
renascimento para se libertar de toda essa ilusão, de todo esse apego a Deus,
apego à doutrina e apego às ideias mais elevadas -- os valores mais nobres.

Todas estas coisas continuam a ser muito boas e louváveis. Mas é através do
apego que sofremos, porque se nos apegamos a qualquer percepção, não estamos a
perceber a verdade. Estamos apenas a apegar-nos a um símbolo e a agarrar-nos
ao símbolo como se fosse a realidade. Se eu dissesse: «Vês o Buda sentado ali em
cima no altar? Esse é o verdadeiro Buda. Aquele é o Buda.» Você pensaria: «O
Ajahn Sumedho realmente enlouqueceu.» No entanto, ainda podemos apegar-nos a
uma percepção do Buda como sendo o Buda. Não podemos? Não estamos na fase em que
vamos acreditar que a enorme estátua é o Buda, mas podemos estar muito apegados
a uma visão sobre o Buda. E pode ser uma visão muito bonita também. Tal como
aquela Buddha-rupa, é uma Buddha-rupa muito bonita, não é? Gosto desse Buda, é
muito bonito. Isso não significa que tenhamos de nos livrar dele, porque as
Buddha-rupas não iludem. O que é perigoso é o nosso apego a uma percepção -- do
eu, dos outros, de Buda, de Deus, da Unidade ou da totalidade.

Quando conseguimos realmente libertar a nossa mente do apego, então todos estes
ângulos específicos são válidos. Não estamos a condenar o monismo como algo
errado. Mas o apego à doutrina monista é limitador e cega-te. Tal como o apego
ao não-dualismo. O propósito do não-dualismo é, na verdade, uma indicação
tremenda do apego. Mas se for apenas um não-dualista filosófico, então pode
ficar apegado a uma atitude de aniquilação.

Não te estou a pedir que te apegues a uma posição de não-dualismo, mas pedi-te
para não tentares inspirar a tua mente ou ler sobre ensinamentos monistas
inspiradores e várias outras religiões durante este retiro, porque, para
aprenderes realmente a usar esta ferramenta específica, tens de a usar e
observar os resultados. E isso pode parecer bastante árido. Mas temos de largar
essa necessidade de inspiração até ao ponto do desespero. Temos de aprender a
aceitar esse vazio, esse silêncio, a cessação, a solidão, a falta de calor, e
não pedir benevolência e bondade. Temos de nos abrir ao silêncio e
contemplá-lo, aprendendo com ele em vez de fugir dele para procurar uma mãe
calorosa ou um pai seguro.

Então, uma forma de descrever esta Vida Santa é o amadurecimento de um ser
individual até àquela maturidade em que já não nos demoramos no calor da
adolescência ou da infância, nem nos prazeres do mundo.

